Information on the locations of refugee camps and centers is drawn from \citet{generalitat_de_catalunya_relacio_2020}. The information in the dataset comes from two original sources in the regional archives: \textit{Dossier on the International Conference for Aid to Spanish Refugees} (Paris, July 15–16, 1939), prepared by the Bureau d'Information (ANC1-511-T-332), and \textit{List of Cities and Towns in France with Camps or Shelters for Spanish Refugees} (ANC1-511-T-349).

%resistance data
To examine the effects on political outcomes, we draw on city-level electoral data from legislative elections compiled by \citet{piketty_histoire_2023}. We restrict the analysis to Metropolitan France and to the years 1919–1956, since 1958 marks the beginning of the Fifth Republic. In order to measure the participation into the French Resistance during the Second World War, we use public available records from the French Ministry of Armies \citep{ministry_of_the_armed_forces_memoire_2022}. Namely, we use records of official recognitions of resistance activities (\textit{Titres, homologations et services pour faits de résistance}), and recipients of resistance medals (\textit{Médaillés de la résistance}). These nominative registers provide information on age, nationality, year and place of birth at the municipality level. We locate involvement in the resistance using place of birth, considering the bad quality of the place of recruitment or death in the dataset. 

%We explore as an outcome also the effect on collaboration with the Nazi Occupation using data on the number of collaborators at the city level from \citet{cage_heroes_2023}.

We use several sources of data to check for the exogeneity of the dispersion policy. Unfortunately, no direct measure of the foreign-born population exists at the municipal level for the period. To approximate pre-existing immigrant presence, we construct a proxy based on national death records from 1970 to 2025 \citep{insee_deces_2025}. From these records, we identify French individuals with Spanish, Italian, or Portuguese surnames and recover their municipality of birth. Using this information, we reconstruct the pre-1939 shares of residents with foreign-origin surnames for each municipality. To identify which surnames are of Spanish origin, we count the frequency of surnames within the population born in Spain and use this to generate a list of common Spanish surnames. Housing availability is proxied using a measure of buildings constructed by 1939 \citep{insee1968logements} and the number of associations is drawn from the legal registry of organisations \citep{dila_journalofficiel_portal}. Finally, we include variables capturing geographic and economic structure—altitude, precipitation, and mining activity \citep{pauling2006,camino2021}. 


%To explore whether the camp locations predict the settlement of the refugees, we use the census of 1962, which is the earliest census available \citep{insee_1962}. To measure the intensity of collective actions after WWII, we rely on various archival sources that enable us to compile an inventory of all political activities carried out by Spanish refugees between 1945 and the fall of the dictatorship in Spain in 1977. We mainly gather information on the two main Spanish unions at the time of the Civil War, the \textit{Unión General de Trabajadores} (General Union of Workers, UGT) and the \textit{Confederación Nacional del Trabajo} (National Confederation of Labor, CNT). For both organizations, we are able to identify local branches of these unions and track their years of activity over time. 

%For the UGT, we use the documentations from the \citep{FFLC_ugt}. More precisely, we use an inventory of all UGT activity outside of Spain during the time of the dictatorship. This inventory gives a list of all municipalities that have had a local section of the UGT, with its date of creation. For the CNT, we use data collected from 2 sources. The first one is the \citep{FAL_cnt}, from which we use the report of general assembly of the movement in exile in 1948. This document enables us to have a precise list of all municipalities in France were Spanish refugees settled down and wished to organize politically into one new local section of the CNT after WWII. Unfortunately, since this does not allow us to track the activity of sections through time, we complete this first source with documentations from the the archives of the \citep{CDHS_cnt}. From this source, we use a large set of documents, especially meeting reports from local sections in exile. This allows us to determine the years of activity of all local sections. Using the same reports, we are able to infer the number of persons affiliated to each local sections using the subscription of news members.

